\section{Thảo luận}
\label{sec:discussion}

\paragraph{Kiến trúc so với huấn luyện.} Câu hỏi ban đầu là ưu thế của DLinear, NLinear đến từ kiến trúc hay từ cách huấn luyện. Ở nghiệm tối ưu, câu trả lời khác nhau giữa hai biến thể. Với DLinear, phần kiến trúc gần như không còn: DLinear không phạt có cùng tập ánh xạ như Linear, và ở $\lambda^*$ hai mô hình chênh không quá \nmBThreeDLinMinusLinAllAbsMax{} trên test (\secref{sec:rq1}). Các khác biệt lớn giữa DLinear và Linear trong bảng công bố, ví dụ \nmCThreeETThOneHSevenTwoZeroPubD{} ở ETTh1 $H = 720$, vì vậy phải đến từ quá trình huấn luyện. Với NLinear, ràng buộc tổng hàng là một khác biệt kiến trúc thật, và nó có lợi rõ trên ETTh2 ngay cả ở nghiệm tối ưu. Khi huấn luyện bằng SGD với lr gốc, DLinear tốt hơn Linear có ý nghĩa ở ETTm1, và lợi thế này biến mất khi lr được tinh chỉnh riêng cho từng mô hình. Trên ETTm1, Linear chọn lr từ $\times 8$ trở lên ở \nmTunedETTmOneLinearChoseGeEightCount/4 horizon, còn DLinear chọn không quá $\times 2$ ở \nmTunedETTmOneDLinearChoseLeTwoCount/4. Điều này phù hợp với \emph{giả thuyết} rằng DLinear có learning rate hiệu dụng lớn hơn Linear ở cùng lr danh nghĩa. Chúng tôi chưa kiểm giả thuyết này.

\paragraph{Outlier ở đích so với outlier ở đầu vào.} Lý thuyết hồi quy bền vững bảo vệ khỏi outlier trong biến đích. Kết quả khi chèn chỉ vào đích khớp với lý thuyết: MAE và Huber gần như không đổi, còn MSE suy giảm theo độ lớn outlier. Khi outlier nằm trong chuỗi, cùng một điểm bất thường xuất hiện ở cả đầu vào của nhiều cửa sổ. Hàm mất mát chỉ hạ trọng số các phần dư lớn, nên không chặn được ảnh hưởng của các hàng đầu vào bất thường. Ở tỷ lệ cao và độ lớn lớn, MAE và Huber suy giảm nhanh hơn MSE.

\paragraph{Giả thuyết nhiễu đầu vào (chưa kiểm).} Ở tỷ lệ $5\%$ và $10\%$ (chèn vào chuỗi), val chọn $\lambda^* = 0$ ở mọi đường $\lambda$: \nmOutLamZeroMseMaeCount/\nmOutLamZeroMseMaeTotalCount{} trường hợp (đường, metric) với MSE và MAE, và cả \nmOutLamZeroCount{} trường hợp với Huber. Một cách giải thích khả dĩ: outlier trong đầu vào tác động như nhiễu cộng vào $x$. Với hàm mất mát bình phương, huấn luyện với nhiễu đầu vào tương đương, ở bậc thấp nhất, với chính quy hóa Tikhonov~\cite{bishop1995training}. Nếu vậy, dữ liệu có outlier đã tự mang một lượng ridge ngầm, nên không cần phạt tường minh nữa, và MSE được ``chính quy hóa sẵn'' theo cách mà MAE không có. Đây là \emph{giả thuyết}. Chúng tôi chưa thiết kế phép kiểm nào cho nó, chẳng hạn so nghiệm MSE trên dữ liệu có outlier trong $x$ với ridge trên dữ liệu sạch ở một $\lambda$ tương ứng.

\paragraph{Nhiễu theo seed trong các bảng công bố.} Bảng của bài gốc báo một lần chạy cho mỗi cấu hình. Ở một số ô, độ lệch chuẩn theo seed của SGD cùng cỡ với chênh lệch giữa các mô hình mà bài báo cáo (\tabref{tab:b1}). Ví dụ ở ETTh2, $H = 720$, độ lệch chuẩn là \nmBOneETThTwoHSevenTwoZeroLinearSd{} (Linear) và \nmBOneETThTwoHSevenTwoZeroDLinearSd{} (DLinear), trong khi chênh trong bài là \nmCThreeETThTwoHSevenTwoZeroPubD. So sánh giữa các mô hình tuyến tính vì vậy cần nhiều seed và khoảng tin cậy. Hoặc, như ở đây, so tại nghiệm tối ưu, nơi không còn nhiễu theo seed.
